\begin{table*}
\small
\centering
\resizebox{\textwidth}{!}{\begin{tabular*}{\textwidth}{l@{\extracolsep{\fill}}cccccccc}
\toprule
\textbf{Model}& \textbf{Overall} & \textbf{Multi-hop}   & \textbf{Qualifier} & \textbf{Comparison} & \textbf{Logical}    & \textbf{Count}  & \textbf{Verify} & \textbf{Zero-shot}   \\ 
\midrule
KVMemNet &16.61  &16.50 &18.47 &1.17 &14.99& 27.31 &54.70& 0.06   \\
EmbedKGQA &28.36 &26.41 &25.20 &11.93 &23.95 &32.88 &61.05 &0.06 \\
RGCN &35.07 &34.00 &27.61 &30.03 &35.85 &41.91 &65.88 &0.00 \\
\midrule
RNN SPARQL &41.98& 36.01 &19.04& 66.98& 37.74 &50.26 &58.84& 26.08   \\
BART SPARQL &89.68 & 88.49 &83.09 &96.12 &88.67 &85.78 &92.33& 87.88 \\
\midrule
ChatGPT & 24.96 & 24.22 & 26.37 & 39.15 & 25.51 & 10.76 & 54.70 & 15.67 \\
Davinci-003 & 31.02 & 29.58 & 31.58 & 49.8& 29.62 & 16.70 & 65.54 & 21.83 \\
GPT-4 & 37.43 & 34.82 & 37.15 & 55.75 & 36.81 & 15.27 & 72.93 & 27.28  \\
\midrule 
Ours & \textbf{92.15} & \textbf{91.03} & \textbf{87.90} & \textbf{96.32} & \textbf{91.28 }& \textbf{88.21} & \textbf{92.86} & \textbf{91.40} \\
\bottomrule
\end{tabular*}}
\caption{The accuracy on the test set of KQA Pro, which is based on Wikidata KG. The results of Davinci-002, GPT-4, and ChatGPT are evaluated by us and the results of other baselines are copied from~\citet{KQA-Pro}.}
\label{tab:wiki-res}
\end{table*}

\section{Experiment}
\label{experiment}
\subsection{Experimental Setup}
\label{sec:app_experiment}
We select four commonly-used KGQA datasets as in-domain datasets, \ie \textit{WebQSP}, \textit{CWQ}, and \textit{GrailQA}, which are based on Freebase, and \textit{KQA Pro}, which is based on Wikidata. And we select three ODQA datasets as out-of-domain datasets, \ie \textit{WQ}, \textit{NQ}, and \textit{TQ}.
Further, we consider three types of baseline methods, \ie \textit{subgraph-based reasoning}, \textit{LM-based seq2seq generation}, and \textit{LLM-based methods} for comparison on in-domain datasets, and \textit{Fine-tune based} and \textit{LLM-based} methods for out-of-domain datasets.
We show the details of the above datasets, baselines, evaluation protocol, and implementation in Appendix~\ref{sec:app-experiment}.


\begin{table}[t]
    % \renewcommand\arraystretch{1.1}
    % \setlength\tabcolsep{5.5pt}
	\centering
	\small
            \begin{tabular}{l  c  c c }
			\toprule
			\textbf{Models} & \textbf{NQ-Wiki} & \textbf{TQ-Wiki} & \textbf{WQ-Freebase} \\
			% \cmidrule{2-4}
			% & Hits@1 & F1 & Hits@1 \\
			\midrule
			T5-Base & 30.94 & 27.63 & 24.06 \\
			T5-Large & 31.21 & 29.40 &  {24.70} \\
			BART-Base & 29.47 & 25.43 & 21.95 \\
			BART-Large	& 32.60 & 33.05 & 26.33 \\
                \midrule
                Davinci-003  & 51.94 & 88.57 & 23.81 \\
			ChatGPT  & \textbf{57.49} & \textbf{88.68}  & 23.23 \\
                \midrule
                Ours & 33.00 & 35.89 & \textbf{28.90} \\
    		\bottomrule
		\end{tabular}
        \caption{The results on the subsets of the dev sets from the out-of-domain ODQA datasets.}%Bold font denotes the best performance.}% We reproduce the performance of all baseline models. Bold font denotes the best performance.}
	\label{tab:out-domain}
\end{table}

\subsection{Main Results}
\paragraph{Results on In-domain Datasets.}
Table~\ref{tab:fb-res} and Table~\ref{tab:wiki-res} show the results on in-domain datasets based on Freebase and Wikidata, respectively. 
First, LM-based seq2seq generation methods can achieve better F1 score compared to the subgraph-based reasoning methods on the WebQSP and KQA Pro. It indicates that the SPARQL query generated by the LM can obtain a more complete answer set, and the structured query can better support some complex operations~(\eg maximum, count) than the traditional subgraph-based reasoning methods. 
Second, although LLMs are powerful, directly using Davinci-003, ChatGPT, and even GPT-4 still has a large performance gap compared with the best fine-tuned methods in WebQSP, GrailQA, and KQA Pro, indicating the difficulty of answering complex questions solely by LLMs. 

Finally, our KG-Agent is substantially better than all other competitive baselines in all datasets after instructing tuning on the mixed data. With the mutual augmentation between different datasets, our approach achieves 1.7\%, 7.5\%, and 2.7\% improvements of F1 on WebQSP, CWQ, and Grailqa, respectively. 
Benefiting from the autonomous reasoning mechnism, our approach can perform reasoning on the two KGs and obtain consistent improvement on all datasets.



\paragraph{Results on Out-of-domain Datasets.}
After instruction tuning, we directly evaluate the zero-shot performance of our KG-Agent on the out-of-domain datasets. As shown in Table~\ref{tab:out-domain}, although fine-tuned with full data, the small pre-trained language models~(\eg T5 and BART) can not effectively answer these factual questions. Owing to the large-scale parameters, Davinci-003 and ChatGPT performs well on NQ and TQ, which are constructed based on Wikipedia, the corpus that they may have been pre-trained on.
However, they perform not well on WQ, which is constructed based on Freebase KG.
%, which is relatively scarce. 
In contrast, our KG-Agent only needs to learn how to interact with KG instead of memorizing the specific knowledge. Thus, it can utilize the external KG in zero-shot setting, and achieve consistent improvement compared to fine-tuned pre-trained language models.
\begin{table}[t]
    \renewcommand\arraystretch{1}
    \setlength\tabcolsep{4.3pt}
	\centering
	\small
            \begin{tabular}{l  c  c c }
			\toprule
			\textbf{Models} & \textbf{MQA-1hop} & \textbf{MQA-2hop} & \textbf{MQA-3hop} \\
			% \cmidrule{2-4}
			% & Hits@1 & F1 & Hits@1 \\
			\midrule
			GraftNet & 82.5    & -  &  -  \\
			EmbedKGQA & 92.0    & 40.7  &  34.6  \\
			  NSM & 94.8  &   97.0    & 91.0 \\
                TransferNet & 96.5    & 97.5  &  90.1 \\
                \midrule
			ChatGPT & 61.9 & 31.0 & 43.2\\
			StructGPT  & 94.2 & 93.9 & 80.2  \\
                \midrule
                Ours & \textbf{97.1} & \textbf{98.0} & \textbf{92.1} \\
    		\bottomrule
		\end{tabular}
        \caption{The results on the three subsets of MetaQA. We copy the results of baselines from~\citet{Jiang-2023-arxiv-StructGPT}.}% Bold font denotes the best performance.}
	\label{tab:res-metaqa}
\end{table}
\subsection{Further Analysis}

\paragraph{Transfer to Domain-specific KG.}
To evaluate the transferability of our approach on other KGs, we test our KG-Agent on the MetaQA dataset which is based on a movie domain KG. 
Following existing work~\cite{NSM, Jiang-2023-arxiv-StructGPT}, we show the one-shot results on the test set in Table~\ref{tab:res-metaqa}. ChatGPT performs not well when directly answering these domain-specific questions, where the performance drops 45\% absolutely on the MQA-3hop subset compared to the supervised fine-tuned TransferNet model. After equipping the LLM with the KG, StructGPT can greatly outperform ChatGPT with about 37\% improvement. In contrast, our KG-Agent can obtain consistent performance improvement compared to the competitive supervised fine-tuning baselines on all subsets. It indicates that the agent indeed learns the general ability about reasoning on KG, which can be efficiently transferred to other KGs.

\paragraph{Effect of Instruction Amount.}
We explore how the amount of instructions affects the performance of KG-Agent and show the results in Figure~\ref{fig:amount}. With a constant sampling proportion, we scale the total amount from 2k to 64k in an exponential way and evaluate the F1 and Hist@1 scores on WebQSP and CWQ datasets. As we can see, the performance increases with more instruction tuning data, and eventually reaches a stable state, which indicates the importance of data amount. At the same time, with the data amount increasing from 16k to 64k, the KG-Agent doesn't obtain a remarkable performance improvement. We think this is relevant to the variety of our instruction tuning data, which is illustrated in existing work~\cite{Chung-arxiv-2022-Scaling, Aribandi-ICLR-2022-ExT5}. Therefore, we will construct more diverse samples in the future to further boost the performance.

\begin{table}[t]
	\centering
	\small
    \begin{tabular}{c | c c c | c}
        \toprule
          Proportion& WebQSP & CWQ & GrailQA & Average\\
        \midrule
            1:10:5 & 80.0& 69.8& 86.1 & \textbf{78.6}  \\
            \underline{2}:10:5 & 81.2& 68.7& 83.3 & 77.8 \\
            1:\underline{20}:5 & 78.9.& 73.6& 78.8 & 77.1 \\
            1:10:\underline{10} & 80.8& 66.9& 84.3 & 77.3 \\
        \bottomrule
    \end{tabular}
        \caption{The F1 scores on three in-domain datasets after instruction tuning under different sampling proportions. We highlight the changed proportion with an underline. }
        \label{tab:ablation}
\end{table}

\paragraph{Effect of Tuning Data Proportion.}
Our experiment finds that only sampling 10K samples from existing datasets is enough for backbone LLM to learn the autonomous decision making capability.
Here, we conduct a further ablation study to explore the impact of sampling proportion on the agent's performance when keeping the total amount of instruction tuning data constant.
Specifically, we evaluate the agent performance of WebQSP, CWQ, and GrailQA when doubling the proportion of one dataset while maintaining the other two dataset proportions. We show the results in Table~\ref{tab:ablation}. We can see that as the sampling proportion of a certain dataset increases, the agent performance on it consistently improves. However, for the average performance on all three datasets, all variants are lower than our selected proportion, indicating that the proportion we chose is suitable for the LLM to balance and master more comprehensive and general abilities.